% titel: Federpendel im Wasser
% skript: Kap. Schwingungen

% Kugel an Feder im Wasserbad: statische Dehnung (Auftrieb), Stokes-Reibung,
% gedämpfte Schwingung, Periodizitätsbedingung, Energiedissipation,
% getriebener Oszillator.

Eine Kugel hängt an einer Feder und schwingt ganz in Wasser
(Auftrieb + Stokes-Reibung berücksichtigt; die Ruhelage ist der neue
Gleichgewichtspunkt um die Dehnung $x_R$).

\begin{werkzeugkasten}[Werkzeuge]
\textbf{Statisches Gleichgewicht} an der Feder (Kraft = Gewicht minus
Auftrieb), Kugelvolumen $V=\frac43\pi R^3$:
\begin{equation}
  D\,x_R=(m g-F_A),\qquad F_A=\rho_{\mathrm W}V g\;,
  \quad m=\rho_{\mathrm K}V\;.
\end{equation}

\textbf{Stoke'sche Reibung} (laminare Kugelumströmung):
\begin{equation}
  F=-\alpha\, v\;,\qquad \alpha=6\pi\eta R\;.
\end{equation}

\textbf{Gedämpfte Schwingung} um die Ruhelage ($\ddot x$ relativ zur Ruhelage):
\begin{equation}
  \ddot x+2\beta\,\dot x+\omega_0^2 x=0,\qquad
  \omega_0=\sqrt{\tfrac{D}{m}},\quad
  \beta=\frac{\alpha}{2m}=\frac{3\pi\eta R}{m}\;.
\end{equation}
\textbf{Periodisch} (unterkritisch) genau dann, wenn $\beta<\omega_0$;
überkritisch ($\beta>\omega_0$) gibt es keine Schwingung.

\textbf{Energie} bei schwacher Dämpfung: $E(t)\approx E_0\,e^{-2\beta t}$;
Verlust pro Periode $\approx 2\beta T$.

\textbf{Getriebener Oszillator} mit Bewegung $z(t)$ der Aufhängung:
\begin{equation}
  \ddot x+2\beta\,\dot x+\omega_0^2 x=\omega_0^2 z(t)\;.
\end{equation}
Amplitudenverhältnis und Phase bei Kreisfrequenz $\omega$:
\begin{equation}
  \frac{x_0}{z_0}=\frac{\omega_0^2}{\sqrt{(\omega_0^2-\omega^2)^2+(2\beta\omega)^2}},
  \qquad \tan\delta=\frac{2\beta\omega}{\omega_0^2-\omega^2}\;.
\end{equation}

\textbf{Einheiten:} $\alpha$ in \unit{\kilogram\per\second},
$\eta$ in \unit{\pascal\second}, $\beta,\omega$ in \unit{\per\second}.
\end{werkzeugkasten}

\begin{werkzeugkasten}[Vorgehen]
\begin{itemize}
  \item Statik: Dehnung aus Kraftgleichgewicht --- Auftrieb nicht vergessen.
  \item Reibungskoeffizient bei Kugeln laminar direkt $\alpha=6\pi\eta R$.
  \item Bewegungsgleichung um die (neue) Ruhelage schreiben, dann $\omega_0$
       und $\beta$ ablesen.
  \item Periodizität: $\beta<\omega_0$; dabei $m$ als Funktion von $R$ einsetzen.
  \item Energiefrage: $E\propto e^{-2\beta t}$; im Ansatz für das getriebene
       Pendel die Aufhängung $z(t)$ als Inhomogenität.
\end{itemize}
\end{werkzeugkasten}

\begin{achtung}[Typische Fallen]
\begin{itemize}
  \item Auftrieb nutzt die \emph{Wasser}-Dichte, nicht die Kugeldichte.
  \item $m$ der Kugel hängt von $R^3$ ab --- bei der Periodizitätsbedingung
       nicht als Konstante behandeln.
  \item Die Ruhelage unter Wasser (um die gedehnte Feder) ist der Ursprung;
       die Schwingung ist \emph{relativ} zu ihr.
  \item Phasendifferenz: bei Resonanz läuft die Auslenkung der Anregung um
       $\pi/2$ nach.
\end{itemize}
\end{achtung}
